\chapter{Entscheiden für viele}

\section{Zwei Programme, dieselbe Summe}

\begin{merksatz}
\boxtitel{Der Fall}
Zwei Programme kosten gleich viel und bringen zusammen je $50$ gewonnene Lebensjahre bei $100$ Patienten.

\begin{center}
{\small
\renewcommand{\arraystretch}{1.5}
\begin{tabular}{@{}P{2.4cm}P{3.0cm}P{2.4cm}@{}}
\toprule
\textbf{Programm} & \textbf{Verteilung} & \textbf{Summe}\\
\midrule
$A$ & jeder $+0{,}5$ Jahre & $50$\\
$B$ & $25$ Patienten $+2$, die übrigen nichts & $50$\\
\bottomrule
\end{tabular}}
\end{center}

\begin{center}
\begin{minipage}{0.9\textwidth}
\textbf{Die Summe entscheidet nicht --- sie ist gleich.}
\end{minipage}
\end{center}

\begin{center}
{\small
\renewcommand{\arraystretch}{1.48}
\begin{tabular}{@{}P{2.6cm}P{2.2cm}P{2.2cm}@{}}
\toprule
\textbf{Maß} & $A$ & $B$\\
\midrule
Summe & $50$ & $50$\\
Gini & $0{,}00$ & $\mathbf{0{,}75}$\\
Summe der Wurzeln & $\mathbf{70{,}7}$ & $35{,}4$\\
\bottomrule
\end{tabular}}
\end{center}

\begin{center}
\begin{minipage}{0.88\textwidth}
\kursiv{\textbf{Die dritte Zeile ist eine Bewertung, keine Messung}: Wer die Wurzel des Gewinns summiert, bevorzugt die gleichmäßige Verteilung --- genau wie eine gekrümmte Nutzenfunktion in Kapitel~2 den sicheren Betrag bevorzugt. Die Krümmung drückt dort Risikoscheu aus, hier Ungleichheitsscheu. Wer nur Summen vergleicht, hat sich für die gerade Linie entschieden.}
\end{minipage}
\end{center}
\end{merksatz}

\section{Vier Einwände gegen die Summe}

\begin{merksatz}
\boxtitel{Was das QALY nicht sieht}
\begin{center}
{\small
\renewcommand{\arraystretch}{1.55}
\begin{tabular}{@{}P{2.6cm}P{6.0cm}@{}}
\toprule
\textbf{Verteilung} & Wer wenig hat, zählt genauso wie wer viel hat.\\
\textbf{Behinderung} & Wer bereits eingeschränkt lebt, kann durch dieselbe Behandlung weniger QALY gewinnen --- und wird deshalb nachrangig.\\
\textbf{Alter} & Wer jünger ist, gewinnt mehr Jahre. Die Rechnung bevorzugt Junge, ohne es zu sagen.\\
\textbf{Rettungspflicht} & Ein identifizierbarer Mensch in Not wiegt in der Wahrnehmung schwerer als viele namenlose.\\
\bottomrule
\end{tabular}}
\end{center}

\schritt{Zur zweiten Zeile, gerechnet}
Eine Behandlung verlängert das Leben um fünf Jahre.
\begin{center}
{\small
\renewcommand{\arraystretch}{1.48}
\begin{tabular}{@{}P{2.8cm}P{2.2cm}P{2.4cm}@{}}
\toprule
\textbf{Patient} & \textbf{Nutzen/Jahr} & \textbf{Gewinn}\\
\midrule
ohne Einschränkung & $1{,}0$ & $5{,}0$ QALY\\
im Rollstuhl & $0{,}7$ & $3{,}5$ QALY\\
\bottomrule
\end{tabular}}
\end{center}

\begin{center}
\begin{minipage}{0.9\textwidth}
\textbf{Dieselbe Behandlung, dieselben fünf Jahre --- und $30$ Prozent weniger Anspruch.}
\end{minipage}
\end{center}

\begin{center}
\begin{minipage}{0.88\textwidth}
\kursiv{\textbf{Das ist der schwerste Einwand gegen die Verwendung des QALY zur Priorisierung.} Er trifft nicht die Rechnung, sondern ihre Verwendung: Als Maß für den Gewinn \kursiv{eines} Patienten ist das QALY sinnvoll; als Maß dafür, wessen Behandlung Vorrang hat, benachteiligt es systematisch die bereits Benachteiligten. Einige Länder haben deshalb Zusatzgewichte eingeführt oder das Maß ganz verboten.}
\end{minipage}
\end{center}
\end{merksatz}

\section{Der Bruchfall}

\begin{warnung}
\boxtitelr{Was passiert, wenn man nicht rechnet}
\begin{center}
{\small
\renewcommand{\arraystretch}{1.52}
\begin{tabular}{@{}P{2.6cm}P{6.0cm}@{}}
\toprule
\textbf{Verdeckte Rationierung} & Wer die Abwägung nicht ausspricht, trifft sie trotzdem --- über Wartezeiten, Zuzahlungen, Zugangshürden.\\
\textbf{Ungleichheit} & Die verdeckte Form trifft regelmäßig die Schwächeren stärker als eine offene Regel.\\
\textbf{Unprüfbarkeit} & Eine unausgesprochene Regel lässt sich nicht kritisieren und nicht verbessern.\\
\bottomrule
\end{tabular}}
\end{center}

\begin{center}
\begin{minipage}{0.88\textwidth}
\kursiv{\textbf{Das ist die Verteidigung der Rechnung, und sie ist stark}: Der Vergleich ist nicht \enquote{Rechnung gegen ärztliches Urteil}, sondern \enquote{offene Regel gegen verdeckte}. Zugleich bleibt der Einwand des vorigen Abschnitts bestehen. Beides zusammen ergibt die vernünftige Haltung: rechnen, das Ergebnis aussprechen --- und die Abweichungen von der Summe ebenfalls begründen, statt sie zu verschweigen.}
\end{minipage}
\end{center}
\end{warnung}

\section{Übungsaufgaben}

\begin{aufgabe}
\boxtitel{Aufgabe 14.1 --- Gewichtete Bewertung}
Bewerten Sie die Programme $A$ und $B$ mit $\sum \sqrt{x_i}$ und mit $\sum x_i$. Welches Gewicht hat die Ungleichheitsscheu?

\loesung
\schritt{Schritt 1 --- Die Summen}
\[ A: \ 100\cdot0{,}5 = 50, \qquad B: \ 25\cdot2 = 50 \]

\schritt{Schritt 2 --- Die Wurzelsummen}
\[ A: \ 100\cdot\sqrt{0{,}5} = 100\cdot0{,}7071 = 70{,}71 \]
\[ B: \ 25\cdot\sqrt{2} = 25\cdot1{,}4142 = 35{,}36 \]

\schritt{Schritt 3 --- Der Umschlagpunkt}
Wie viel müsste $B$ je Begünstigtem bringen, um gleichzuziehen?
\[ 25\sqrt{g} = 70{,}71 \quad\Longrightarrow\quad \sqrt{g} = 2{,}828 \quad\Longrightarrow\quad g = 8{,}0 \]

\begin{center}
\begin{minipage}{0.9\textwidth}
\textbf{$B$ müsste jedem Begünstigten acht statt zwei Jahre bringen --- also die vierfache Gesamtsumme.}
\end{minipage}
\end{center}

\begin{center}
\begin{minipage}{0.88\textwidth}
\kursiv{\textbf{Die Wurzel ist eine sehr starke Ungleichheitsscheu.} Man kann sie abschwächen ($x^{0{,}8}$) oder verschärfen ($\ln x$); jede Wahl ist eine Setzung, und sie sollte ausgesprochen werden. Wichtiger als die Wahl selbst ist, dass sie einmal getroffen und dann auf alle Programme gleich angewandt wird --- sonst entscheidet wieder der Einzelfall.}
\end{minipage}
\end{center}
\end{aufgabe}

\begin{aufgabe}
\boxtitel{Aufgabe 14.2 --- Das Alter in der Rechnung}
Zwei Patienten, $30$ und $70$ Jahre alt, brauchen dieselbe Behandlung; es gibt nur eine. Wie entscheiden verschiedene Maßstäbe?

\loesung
\schritt{Schritt 1 --- Gewonnene Jahre}
\begin{center}
{\small
\renewcommand{\arraystretch}{1.48}
\begin{tabular}{@{}P{2.0cm}P{2.4cm}P{2.6cm}@{}}
\toprule
\textbf{Alter} & \textbf{Restleben} & \textbf{Gewinn}\\
\midrule
$30$ & $50$ Jahre & $\mathbf{50}$\\
$70$ & $15$ Jahre & $15$\\
\bottomrule
\end{tabular}}
\end{center}

\schritt{Schritt 2 --- Drei Maßstäbe}
\begin{center}
{\small
\renewcommand{\arraystretch}{1.5}
\begin{tabular}{@{}P{2.8cm}P{2.4cm}P{2.6cm}@{}}
\toprule
\textbf{Maßstab} & \textbf{Wahl} & \textbf{Begründung}\\
\midrule
gewonnene Jahre & $30$-Jähriger & mehr Nutzen\\
gleiche Chance & Los & gleiches Recht\\
faire Lebensspanne & $30$-Jähriger & hatte weniger\\
\bottomrule
\end{tabular}}
\end{center}

\begin{center}
\begin{minipage}{0.88\textwidth}
\kursiv{\textbf{Bemerkenswert ist, dass der erste und der dritte Maßstab zum selben Ergebnis kommen --- aus entgegengesetzten Gründen.} Der eine zählt den künftigen Nutzen, der andere das bisher Erhaltene. In der Zuteilung knapper Intensivbetten wurden 2020 beide Begründungen benutzt; die deutschen Empfehlungen entschieden sich ausdrücklich \kursiv{gegen} das Alter als eigenständiges Kriterium und für die kurzfristige Überlebenswahrscheinlichkeit.}
\end{minipage}
\end{center}

\schritt{Schritt 3 --- Warum das kein Rechenproblem ist}
\begin{center}
\begin{minipage}{0.9\textwidth}
\textbf{Alle drei Maßstäbe sind widerspruchsfrei. Die Mathematik sagt nicht, welcher gilt.}
\end{minipage}
\end{center}
\end{aufgabe}

\begin{aufgabe}
\boxtitel{Aufgabe 14.3 --- Der Leitfall, zum letzten Mal}
Frau B.\ soll das neue, teurere Mittel erhalten. Aus Kapitel~10: ICER $45\,004$ Euro je QALY, Schwelle $20\,000$. Was folgt?

\loesung
\schritt{Schritt 1 --- Aus Sicht der Patientin}
\[ +0{,}26 \ \text{QALY} \quad\Longrightarrow\quad \text{das Mittel ist besser} \]

\schritt{Schritt 2 --- Aus Sicht der Gesamtheit}
\[ 11\,656 \ \text{Euro verdrängen} \ \frac{11\,656}{20\,000} = 0{,}58 \ \text{QALY anderswo} \]
\[ \text{Bilanz: } +0{,}26 - 0{,}58 = \mathbf{-0{,}32} \ \text{QALY} \]

\begin{center}
\begin{minipage}{0.9\textwidth}
\textbf{Für Frau B.\ ist es die bessere Behandlung. Für die Bevölkerung ist es ein Verlust.}
\end{minipage}
\end{center}

\schritt{Schritt 3 --- Was folgt daraus \kursiv{nicht}}
\begin{center}
{\small
\renewcommand{\arraystretch}{1.5}
\begin{tabular}{@{}P{2.8cm}P{5.8cm}@{}}
\toprule
\textbf{Nicht:} & Der Arzt soll am Bett die Bevölkerungsrechnung anstellen.\\
\textbf{Sondern:} & Die Regel, welche Mittel zur Verfügung stehen, wird an anderer Stelle gesetzt --- und dort gehört die Rechnung hin.\\
\bottomrule
\end{tabular}}
\end{center}

\begin{center}
\begin{minipage}{0.88\textwidth}
\kursiv{\textbf{Diese Arbeitsteilung ist der einzige Weg, beide Loyalitäten zu wahren}: Der Arzt vertritt seinen Patienten, das Gemeinwesen setzt den Rahmen. Wer beides in einer Person vereint, macht den Arzt zum Sachwalter fremder Interessen am Krankenbett --- und zerstört das Vertrauen, ohne die Verteilungsfrage besser zu lösen. Die dritte Warnung des Vorworts ist damit keine Rechenaufgabe, sondern eine Frage der Zuständigkeit.}
\end{minipage}
\end{center}
\end{aufgabe}

\begin{geschichte}
\boxtitelg{Oregon und die Liste}
Der US-Bundesstaat Oregon stellte 1990 eine Rangliste von $709$ Leistungen auf, um die staatliche Krankenversicherung auf mehr Menschen auszudehnen.

\begin{center}
\begin{minipage}{0.9\textwidth}
\kursiv{Die erste Fassung ordnete rein nach Kosten je QALY --- und setzte die Zahnkappe vor die Blinddarmoperation, weil sie billiger war und oft gebraucht wird. \textbf{Die Liste wurde zurückgezogen und nach ärztlichem Urteil neu geordnet}; erst diese Fassung wurde umgesetzt. Sie gilt als Lehrstück dafür, dass eine Rechnung, die formal richtig ist, an der öffentlichen Beurteilung scheitern kann.}
\end{minipage}
\end{center}

\schritt{Und ein Verbot}
\begin{center}
\begin{minipage}{0.9\textwidth}
\kursiv{Die amerikanische Gesundheitsreform von 2010 untersagte der eingerichteten Forschungsagentur ausdrücklich, Kosten je QALY als Schwelle zu verwenden --- mit der Begründung, das Maß benachteilige Menschen mit Behinderungen und ältere Patienten. \textbf{In Deutschland wird der Nutzen neuer Arzneimittel bewertet, ein ausdrücklicher Schwellenwert je QALY aber nicht genannt}; die Abwägung geschieht in Preisverhandlungen.}
\end{minipage}
\end{center}
\end{geschichte}

\begin{forschung}
\boxtitelp{Verteilung heute}
\begin{center}
{\small
\renewcommand{\arraystretch}{1.45}
\begin{tabular}{@{}P{3.4cm}P{4.8cm}@{}}
\toprule
\textbf{Verteilungsbezogene Analyse} & Wer gewinnt und wer verliert, nicht nur wie viel.\\
\textbf{Gewichte aus Befragungen} & Wie stark die Bevölkerung Ungleichheit scheut.\\
\textbf{Verfahren statt Formel} & Rechtfertigung durch nachvollziehbare Verfahren.\\
\bottomrule
\end{tabular}}
\end{center}

\schritt{Zur dritten Zeile}
\begin{center}
\begin{minipage}{0.88\textwidth}
\kursiv{\Person{Norman Daniels} und \Person{James Sabin} schlugen vor, Priorisierungen nicht durch eine Formel zu rechtfertigen, sondern durch ein Verfahren: öffentliche Begründung, einschlägige Gründe, Berufungsmöglichkeit, Durchsetzung. \textbf{Das verschiebt die Frage von \enquote{welche Rechnung ist richtig} zu \enquote{welches Verfahren ist fair}} --- und entspricht genau dem, was Satz 13.1 nahelegt: Da es die richtige Aggregation nicht gibt, muss die Legitimität aus dem Weg zur Entscheidung kommen.}
\end{minipage}
\end{center}
\end{forschung}

\begin{merksatz}
\boxtitel{Was Sie jetzt können}
\begin{center}
\renewcommand{\arraystretch}{1.5}
\begin{tabular}{@{}P{3.0cm}P{5.6cm}@{}}
\toprule
\textbf{Gleiche Summe} & sehr verschiedene Verteilung\\
\textbf{Ungleichheits\-scheu} & Krümmung wie bei der Risikoscheu\\
\textbf{Vier Einwände} & Verteilung, Behinderung, Alter, Rettung\\
\textbf{Verdeckt} & wer nicht rechnet, teilt trotzdem zu\\
\textbf{Zuständigkeit} & Bett und Rahmen trennen\\
\bottomrule
\end{tabular}
\end{center}

\begin{center}
\begin{minipage}{0.88\textwidth}
\kursiv{Damit endet Teil~V. Das letzte Kapitel zieht Bilanz --- und kehrt ein letztes Mal zu Frau B.\ zurück.}
\end{minipage}
\end{center}
\end{merksatz}
